% Annexe : fiches méthode

\section{Choisir un modèle probabiliste}

\begin{Meth}
Avant tout calcul, poser les quatre questions suivantes :
\begin{enumerate}
    \item Quelle est la grandeur ou le phénomène observé ?
    \item Est-ce un comptage ou une mesure ?
    \item Quelles hypothèses de simplification accepte-t-on ?
    \item Quelle est la question finale : calcul, comparaison, décision ?
\end{enumerate}
\end{Meth}

\section{Fiche réflexe : événements et conditionnement}

\begin{itemize}
    \item Traduire l'énoncé en événements nommés clairement.
    \item Si l'énoncé parle de \og sachant \fg{}, penser à la probabilité conditionnelle.
    \item Si l'énoncé comporte plusieurs cas exclusifs, penser à une partition puis aux probabilités totales.
    \item Si l'on inverse une information, penser a Bayes.
\end{itemize}

\section{Fiche réflexe : lois discrètes}

\begin{center}
\renewcommand{\arraystretch}{1.5}
\begin{tabular}{|m{3cm}|m{4.5cm}|m{4.8cm}|}
\hline
\textbf{Loi} & \textbf{Quand l'utiliser ?} & \textbf{Paramètres et grandeurs utiles} \\
\hline
Bernoulli & Un seul essai, succès ou échec & $p$, moyenne $p$, variance $p(1-p)$ \\
\hline
Binomiale & $n$ essais indépendants, même probabilité de succès & $n$, $p$, moyenne $np$, variance $np(1-p)$ \\
\hline
Géométrique & Rang du premier succès dans des essais indépendants & $p$, moyenne $\dfrac{1}{p}$ \\
\hline
Poisson & Comptage d'événements rares sur une durée ou une zone & $\lambda$, moyenne $\lambda$, variance $\lambda$ \\
\hline
\end{tabular}
\end{center}

\section{Fiche réflexe : lois continues}

\begin{center}
\renewcommand{\arraystretch}{1.5}
\begin{tabular}{|m{3cm}|m{4.5cm}|m{4.8cm}|}
\hline
\textbf{Loi} & \textbf{Quand l'utiliser ?} & \textbf{Paramètres et grandeurs utiles} \\
\hline
Uniforme & Toutes les valeurs d'un intervalle sont équiprobables & $[a,b]$, moyenne $\dfrac{a+b}{2}$ \\
\hline
Exponentielle & Temps d'attente ou durée de vie simplifiée & $\lambda$, moyenne $\dfrac{1}{\lambda}$ \\
\hline
Normale & Mesures influencées par de nombreux petits effets & $\mu$, $\sigma$, centrage-réduction \\
\hline
\end{tabular}
\end{center}

\section{Fiche réflexe : loi normale}

Si $X\sim \mathcal{N}(\mu,\sigma^2)$, on pose
$$
Z=\frac{X-\mu}{\sigma}.
$$

Alors :
\begin{itemize}
    \item $\mathbb{P}(X\leq a)=\mathbb{P}\left(Z\leq \dfrac{a-\mu}{\sigma}\right)$ ;
    \item $\mathbb{P}(a\leq X\leq b)=\mathbb{P}\left(\dfrac{a-\mu}{\sigma}\leq Z \leq \dfrac{b-\mu}{\sigma}\right)$.
\end{itemize}

Valeurs usuelles à connaître :
\begin{itemize}
    \item $1.645$ pour $90\%$ ;
    \item $1.96$ pour $95\%$ ;
    \item $2.576$ pour $99\%$.
\end{itemize}

\section{Fiche réflexe : intervalles de confiance}

\begin{itemize}
    \item IC d'une moyenne :
    $$
    \overline{x} \pm z_{\alpha/2}\frac{\sigma}{\sqrt{n}}.
    $$
    \item IC d'une proportion :
    $$
    \widehat{p} \pm z_{\alpha/2}\sqrt{\frac{\widehat{p}(1-\widehat{p})}{n}}.
    $$
\end{itemize}

\section{Checklist de rédaction}

\begin{itemize}
    \item Nommer la variable ou les événements.
    \item Justifier le choix du modèle.
    \item Écrire la formule avant l'application numérique.
    \item Donner une unité quand il y en a une.
    \item Conclure par une phrase interprétable par un ingénieur.
\end{itemize}

\section{Erreurs classiques à relire avant l'examen}

\begin{itemize}
    \item Confondre $\mathbb{P}(A|B)$ et $\mathbb{P}(B|A)$.
    \item Confondre variable discrète et variable continue.
    \item Oublier le facteur $\sqrt{n}$ dans les calculs statistiques.
    \item Confondre moyenne théorique et moyenne observée.
    \item Écrire une conclusion sans revenir à la question du contexte.
\end{itemize}
